\section{Homework 1}

\Yang{Very nasted, need to revised or just use the typst version.}

In the following, we work over a field $kk$.

\begin{exercise}
    Let $V \subset \bbA^n$ be a closed subvariety of dimension $< n$ with $0 \notin V$ and let $W \subset \mathbb{P}^1 \times \bbA^n$ be the closure of the image of the multiplication map 
    \[
    \left(\bbA^1 \ \left\{0\right\}\right) \times V \to \left(\bbA^1 \ \left\{0\right\}\right) \times \bbA^n , \quad \left(t , v\right) \mapsto \left(t , t v\right).
    \] 
    For the first projection $W \to \mathbb{P}^1$, show that: 
    \begin{itemize}
        \item the fiber over $t = 1$ is $V$; 
        \item the fiber over $t = 0$ is the affine cone of $\overline{V} \cap H$ where $\overline{V} \subset \mathbb{P}^n$ is the projective closure of $V$ and $H \subset \mathbb{P}^n$ is the hyperplane at infinity;
        \item the fiber over $t =\infty$ is empty.
    \end{itemize}
    Under the assumption of $\# kk = \infty$, deduce that $CH_i \left(\bbA^n\right) = 0$ for all $i < n$. 
\end{exercise}

\begin{proof}
    Denote the image of the multiplication map by $U$. 
    Note that $U$ is closed in $\mathbb{G}_m \times \bbA^n$ since it can be defined by $\left\{f\left(x/t\right) : f \in I\left(V\right)\right\}$. 
    Hence the fiber of $W$ over $t = 1$ is equal to the fiber of $U$ over $t = 1$, which is $V$.

    On the affine chart $\bbA^1 \times \bbA^n$ we consider the closure of $U$. 
    For every polynomial $f \in kk\left[X_1 , \ldots , X_n\right]$, denote by $\deg f$ the total degree of $f$. Consider the homogenization $f^h \in kk\left[T , X_1 , \ldots , X_n\right]$ of $f$ with respect to the variable $T$ (i.e., $f^h \left(T , X_1 , \ldots , X_n\right) = T^{\deg f} f\left(\frac{X_1}{T} , \ldots , \frac{X_n}{T}\right)$). Then the closure of $U$ in $\bbA^1 \times \bbA^n$ is defined by the ideal generated by $\left\{f^h \left(T , X_1 , \ldots , X_n\right) : f \in I\left(V\right)\right\}$. Restricting to the fiber over $t = 0$, we get the affine cone of $\overline{V} \cap H$.

For the fiber over $t = \infty$, we can work on the affine chart $\bbA^1 \times \bbA^n$ with coordinates $\left(s , x_1 , \ldots , x_n\right)$ where $s = \frac{1}{t}$. On this chart, the polynomials $f\left(\frac{x}{t}\right) = f\left(s x\right)$ is well-defined. Note that $0 \notin V$, then there exists $f \in I\left(V\right)$ with nonzero constant term. Then the function $f\left(s x\right)$ vanishing on $U$ but is nonzero on the whole fiber over $s = 0$. Hence the fiber of $W$ over $t = \infty$ is empty.

Assume the $\# kk = \infty$. Then there exists a $kk$-point $x \in \bbA^n \ V$. Let $\varphi :\bbA^n arrow \bbA^n$ be an automorphism sending $0$ to $x$. Then $\varphi^{- 1}\left(V\right)$ is a closed subvariety of $\bbA^n$ of dimension $< n$ with $0 \notin \varphi^{- 1}\left(V\right)$. By the above, we have $\left[\varphi^{- 1}\left(V\right)\right] = \left[p^{- 1}\left(0\right)\right]-\left[p^{- 1}\left(\infty\right)\right] \tilde 0$ in $CH_i \left(\bbA^n\right)$ where $p : W arrow \mathbb{P}^1$ is the first projection. Thus $\left[V\right] = \varphi^{- 1 , *} \left[\varphi^{- 1}\left(V\right)\right] \tilde 0$ in $CH_i \left(\bbA^n\right)$. Since $V$ is arbitrary, we get $CH_i \left(\bbA^n\right) = 0$ for all $i < n$. \end{proof}

\begin{remark}
by author of note To show $CH_i \left(\bbA^n\right) = 0$, considering the addition map rather than the multiplication map may be more straightforward. \end{remark}

Let $X \subset \mathbb{P}^2$ be a smooth conic curve. Show that $\deg : CH_0\left(X\right) arrow \mathbb{Z}$ is injective, with image $\mathbb{Z}$ if $X$ has a $kk$-point and $2\mathbb{Z}$ otherwise. (Hint: Use the Riemann-Roch theorem for curves. For the image, as an intermediate step, show that if $X$ has a $kk$-point, then $X \cong \mathbb{P}^1$.) 

\begin{proof}
We have $p_a \left(C\right) = \frac{(\left(d - 1\right)\left(d - 2\right))}{2}$ for a smooth plane curve $C$ of degree $d$. Hence $p_a \left(X\right) = 0$. Let $\alpha \in CH_0\left(X\right)$ of degree $0$. By Riemann-Roch, we have $h^0\left(\alpha\right) - h^0\left(K_X - \alpha\right) = 1$. In particular, $h^0 \left(\alpha\right) \geqslant 1$. Similarly, we have $h^0 \left(-\alpha\right) \geqslant 1$. Hence $\mathcal{O}_X \left(\alpha\right) \cong \mathcal{O}_X$ and $\alpha \tilde 0$ in $CH_0\left(X\right)$.

If $X$ has a $kk$-point, saying $P$, then $\deg \left[P\right] = 1$ and hence $\deg : CH_0\left(X\right) arrow \mathbb{Z}$ is surjective.

If $\deg$ has image $\mathbb{Z}$, then there exists $\alpha \in CH_0\left(X\right)$ with $\deg \alpha = 1$. By Riemann-Roch, we have $h^0 \left(\alpha\right) \geqslant 1$. Hence there exists a $kk$-point $P$ on $X$. In particular, if $X$ has no $kk$-point, then the image of $\deg$ is not $\mathbb{Z}$. Consider the interinterion of $X$ with a line $L$ in $\mathbb{P}^2$. The interinterion $X \cap L$ consists of at most two points counting multiplicity. If $X \cap L$ consists of two distinct points or one point with multiplicity $2$, then $X$ has $kk$-points, a contradiction. Hence $X \cap L$ consists of a single point $P$ whose residue field is a quadratic extension of $kk$. Then $\deg P = 2$ and hence the image of $\deg$ is $2\mathbb{Z}$. \end{proof}

Work out Fulton, Example 2.4.5. Namely, with the notation of the example, show $\left[D\right] = 2\left[l\right]$ and $D \cdot \left[l\right] = \left[P\right]$ and deduce that $\left[l\right]$ is not linear equivalent to any Cartier divisor on $X$. 

\begin{example}
Fulton, Example 2.4.5 Let $\left(x : y : z : t\right)$ be homogeneous on $\mathbb{P}^3$, and let $X$ be the singular cone defined by the equation $z^2 = x y$. Let $D$ be the Cartier divisor on $X$ defined by the equation $x = 0$. Let $l$ be the line $x = z = 0$, $l$ the line $y = z = 0$, $P$ the point $\left(0 : 0 : 0 : 1\right)$. Then $\left[D\right] = 2\left[l\right]$ and $D \cdot \left[l\right] = \left[P\right]$. It follows from Theorem 2.4 that there cannot be a Cartier divisor $D$ on $X$ with $\left[D\right] = \left[l\right]$, either as cycles or as classes in $A_1\left(X\right)$. \end{example}

\begin{proof}

\end{proof}

Work out Fulton, Example 2.5.2. Namely, show (b), (c), (d). (Hint: Use the additivity of Hilbert polynomials. For (b), use Hartshorne, 1.7.4. For (d), use Fulton, A.2.7.) 

\begin{example}
Fulton, Example 2.5.2 Let $X$ be a closed subscheme of $\mathbb{P}^n$ of dimension $\leq k$, and let $\Gamma \left(X\right)$ be the homogeneous coordinate ring of $X$.

 \begin{itemize}
 \item (a) For $t$ sufficiently large, the dimension of the $t^{\text{th}}$ graded piece $\Gamma \left(X\right)_t$ is a polynomial in $t$ of degree $\leq k$ called the Hilbert polynomial (cf. Hartshorne (5) I.7.5). Define $d_k \left(X\right)$ to be the coefficient of $\frac{t^k}{k !}$ in this polynomial.
 \item (b) If $X_1 , \ldots , X_r$ are the irreducible components of $X$, $m_i$ the multiplicity of $X_i$ in $X$, then $d_k \left(X\right) = \sum m_i d_k \left(X_i\right)$.
 \item (c) If $X$ is an irreducible variety, and $H$ is a hypersurface of degree $m$ not containing $X$, then $d_{k - 1} \left(X \cap H\right) = m d_k \left(X\right)$. (Here $X \cap H$ is the scheme-theoretic intersection, so there are exact sequences \[
0 \rightarrow \Gamma(X)_{t - m} \rightarrow \Gamma(X)_t \rightarrow \Gamma(X \cap H)_t \rightarrow 0 \quad \text{for} t \gg 0 .
\]
 \item (d) For any purely $k$-dimensional subscheme $X$ of $\mathbb{P}^n$, \[
d_k \left(X\right) = \deg\left[X\right] = \int_{\mathbb{P}^n} c_1 \left(O\left(1\right)\right)^k \cap \left[X\right].
\]
 \end{itemize}
(For $X$ a subvariety, $c_1 \left(O\left(1\right)\right) \cap \left[X\right] = \left[X \cap H\right]$, where $H$ is a hyperplane not containing $X$.) \end{example}

\begin{proof}
$\mathcal{I} \mathcal{J}$ \end{proof}